Debates about AI writing often treat direct lived experience as a precondition for authentic language. Death is a sharp stress case for that claim: large language models have not experienced death, but human writers also routinely write meaningfully about death without having been dead. We test whether bereavement is therefore a uniquely disqualifying topic for LLM writing, or whether model outputs behave more like human writing about inaccessible experiences than critics assume.

We construct a matched corpus with 12 human grief passages from \aipsy and 24 blinded generations from \gptfive and \claudesonnet. Two independent LLM judges rate each text for authenticity, emotional plausibility, and specificity, and a pretrained \mage detector scores machine-likeness. We then compare human and model texts within bereavement and non-bereavement grief domains.

The human-model authenticity gap is small in both domains. In bereavement, human texts average 4.92/5 authenticity versus 4.75/5 for model texts; in other grief, the corresponding scores are 4.67 and 4.54. Neither contrast is statistically significant, and the source-by-domain interaction is near zero ($+0.042$, $p=0.835$). At the same time, the detector still separates generated from human text with AUROC $=0.826$, while all judge source guesses label every text as human.

This pilot result supports a narrower but important claim: in this setting, death does not function as a uniquely disqualifying topic for LLM writing. Reader-facing authenticity and detector separability are not the same phenomenon.
